\section{Sufficient Conditions for Gradient Estimates}


In this section, we establish \ref{eq:CSGE} by choosing appropriate transition rates and weights, and the resulting bound can be controlled by norms of the Dobrushin influence matrix.

\subsection{Canonical transition rates and weights}
Given a distribution $\mu$ supported on $\Omega=\{\pm 1\}^n$, for any $x \in \Omega$ and $i \in [n]$, we define
\begin{align}\label{eq:canonical-r-w}
    q^\star_i(x) = \frac 12\sqrt{\frac{\mu(x^i)}{\mu(x)}}
        \qquad\text{and}\qquad
    w^\star_i(x) = 2q^\star_i(x) = \sqrt{\frac{\mu(x^i)}{\mu(x)}}.
\end{align}
It is straightforward to show that such transition rates $\{q^\star_i\}_{i \in [n]}$ satisfy the detailed balance condition and hence the associated heat semigroup $(P_t)_{t\ge 0}$ is reversible with respect to $\mu$.
We call \Cref{eq:canonical-r-w} the \emph{canonical transition rates and weights} associated with $\mu$. 

Under the canonical rates and weights, many notions in the $\Gamma$-calculus simplify to a nice form, which we summarize below.
Throughout, we use superscript $\star$ to indicate that the weight functions are canonical, and thereby omit $w^\star$ in the subscript; e.g., $\Gamma^\star = \Gamma_{w^\star}$ denotes the weighted carr\'e du champ operator under $q^\star$ and $w^\star$.
We have the following:
\begin{align*}
    \delta^\star_i f(x) &= \sqrt{\frac{1}{2} w^\star_i(x) q^\star_i(x)} (f(x^i) - f(x)) = q^\star_i(x) (f(x^i) - f(x)); \\
    \grad^\star f &= (\delta^\star_1 f, \dots, \delta^\star_n f); \quad
    L f = \sum_{i=1}^n \delta^\star_i f; \quad
    %\Gamma_{w,i} f(x) = q_i(x)^2 (f(x^i) - f(x))^2 \\
    \Gamma^\star f = \sum_{i=1}^n (\delta^\star_i f)^2.
\end{align*}

We also record two important properties of the canonical transition rates and weights. These properties are essential for establishing \ref{eq:SGE} and \ref{eq:CSGE}.

\begin{fact}[Square property]\label{fact:square}
    For any $i,j \in [n]$ and $x \in \Omega$, we have
    \begin{align*}
        q^\star_i(x) q^\star_j(x^i) = q^\star_j(x) q^\star_i(x^j).
    \end{align*}
\end{fact}

\begin{lemma}
\label{lem:commutator}
    Let $\mu$ be a distribution supported on $\Omega = \{\pm 1\}^n$.
    Let $M: \Omega \to \R^{n \times n}$ be a matrix field defined by, for each $x \in \Omega$,
    \begin{align}\label{eq:local-matrix}
        M_{ij}(x) := 
        \begin{cases}
            \frac{1}2 \tp{\sqrt{\frac{\mu(x^{ij})}{\mu(x^i)}}-\sqrt{\frac{\mu(x^j)}{\mu(x)}}}, & i \neq j;\\
            0, & i = j.
        \end{cases}
    \end{align}
    For any function $f: \Omega \to \R$,
    \begin{align*}
        (L \grad^\star - \grad^\star L) f = \left( \diag\{M\*1\} - M^\top \right) \grad^\star f.
    \end{align*}
\end{lemma}

\begin{proof}
    % \ZC{@Zejia: You need to make the proof more comprehensible. For example, for this proof, first calculate $\delta^\star_j \delta^\star_i - \delta^\star_i \delta^\star_j$, then $L \delta^\star_i - \delta^\star_i L$, and finally $L \grad^\star - \grad^\star L$ (the latter two should follow straightforwardly). Use
    % \begin{align*}
    %     \grad^\star f &= (\delta^\star_1 f, \dots, \delta^\star_n f);\\
    %     L f &= \sum_{i=1}^n \delta^\star_i f;
    % \end{align*}
    % which were mentioned above, to simplify the presentation (doesn't mean shorter proof, but easier to understand). Your proof should rely on the square property, and you should explicitly mention it; this helps the reader understand why the canonical rates and weights work.}
    
    
    Fix $f: \Omega\mapsto \R$, $x\in \Omega$.
    For any $i,j \in [n]$ with $i\neq j$, by definition,
    $$
        \delta^\star_j \delta^\star_i f(x) = q^\star_j(x) (\delta^\star_i f(x^j) - \delta^\star_i f(x)) = q^\star_j(x) q^\star_i(x^j)(f(x^{ij})-f(x^j)) - q^\star_j(x)q^\star_i(x)(f(x^i)-f(x)).
    $$
    \Cref{fact:square} implies that the cancellation of the second derivative in the commutator term $\delta^\star_i\delta^\star_j - \delta^\star_j\delta^\star_i$. 
    We have that,
    \begin{align}
        (\delta^\star_j \delta^\star_i - \delta^\star_i \delta^\star_j) f(x) &= (q_j^\star(x)q_i^\star(x^j) - q_i^\star(x)q_j^\star(x^i)) f(x^{ij})  \nonumber
        \\&-(q_i^\star(x^j) - q_i^\star(x))q_j^\star(x)f(x^j)+(q_j^\star(x^i) - q_j^\star(x))q_i^\star(x)f(x^i) \nonumber
        \\&= M_{ij}(x)q_i^\star(x)f(x^i) - M_{ji}(x)q_j^\star(x)f(x^j), \nonumber
        \\&= M_{ij}(x)\delta^\star_i f(x) - M_{ji}(x)\delta^\star_j f(x), \label{eq:commutator_delta_ij}
    \end{align}
    where the last equality follows from $M_{ij}(x) q_i^\star(x) = M_{ji}(x) q_j^\star(x)$ by \Cref{fact:square}.
    Then the following equality immediately follows from
    \Cref{eq:commutator_delta_ij},
    \begin{align}
        (L\delta^\star_i - \delta^\star_iL) f(x) &= (\sum_{j=1}^n \delta^\star_j\delta^\star_i - \delta^\star_i\delta^\star_j)f(x)
        = \sum_{j=1}^n M_{ij}(x)\delta^\star_i f(x) - \sum_{j=1}^n M_{ji}(x)\delta^\star_j f(x) \nonumber
        \\&= \tp{M(x)\*1}_i\delta^\star_i f(x) - \inner{(M^\top(x))_i}{\nabla^\star f(x)}, \nonumber
    \end{align}
    that is, $(L \grad^\star - \grad^\star L) f = \left( \diag\{M\*1\} - M^\top \right) \grad^\star f$ holds for any $x\in \Omega$.
\end{proof}




\subsection{Sufficient condition for SGE} \label{sec:SGE_boolean_hypercube}

\begin{theorem}[SGE under canonical transition rates and weights]
\label{thm:SGE_boolean_hypercube}
    Let $\mu$ be a distribution on $\Omega = \{\pm 1\}^n$ with full support. Let $M,D: \Omega \to \R^{n \times n}$ be two matrix fields where $M$ is defined by \Cref{eq:local-matrix} and $D$ is defined by, for each $x \in \Omega$,
    \begin{align}\label{eq:def-D}
        D(x) := \diag\left\{\sqrt{\frac{\mu(x)}{\mu(x^i)}}\right\}_{i\in[n]}.
    \end{align}
    Then, under the canonical transition rates and weights, \ref{eq:SGE} holds with constant
    \begin{align*}
        \rho = \min_{x \in \Omega} \;
        \lambda_{\min}\left( D(x) + \diag\{M(x)\*1\} - \frac{1}{2}\left( M(x) + M(x)^\top \right) \right),
    \end{align*}
    where $\lambda_{\min}(A)$ denotes the minimum eigenvalue of a symmetric matrix $A$.
\end{theorem}

\begin{proof}
    By \Cref{prop:equivalence_BE_GE}, it is sufficient to establish the \ref{eq:rBE} condition.
    Specifically, for any function $f: \Omega \to \R$, we aim to show the following pointwise inequality:
    \begin{align*}
        \Gamma^\star_2(f) - \Gamma(\sqrt{\Gamma^\star(f)}) \ge \rho \cdot \Gamma^\star(f),
    \end{align*}
    where we recall that $\Gamma$ is the unweighted carr\'e du champ operator with respect to the canonical rates $q^\star$.
    By \Cref{lem:gamma2-formula}, we have
    \begin{align}
        \Gamma^\star_2(f) - \Gamma(\sqrt{\Gamma^\star(f)}) 
        = \sqrt{\Gamma^\star(f)} L \sqrt{\Gamma^\star(f)} - \Gamma^\star(f,Lf) 
        = \norm{\nabla^\star f}_2 L \norm{\nabla^\star f}_2 - \inner{\nabla^\star f}{\nabla^\star Lf}. \label{eq:first-rBE}
        %&= \tp{ \norm{\nabla^\star f}_2 L \norm{\nabla^\star f}_2 - \inner{\nabla^\star f}{L \nabla^\star f} } + \inner{\nabla^\star f}{(L \nabla^\star - \nabla^\star L) f}.
    \end{align}

    
    Fix $x\in \Omega$. For the first term, we have that
    \begin{align}
        \norm{\nabla^\star f}_2 L \norm{\nabla^\star f}_2 (x)
        &= \norm{\nabla^\star f(x)}_2 \sum_{i=1}^n q^\star_i(x) \tp{\norm{\nabla^\star f(x^i)}_2 - \norm{\nabla^\star f(x)}_2} \nonumber \\
        &= \sum_{i=1}^n q^\star_i(x) \tp{\norm{\nabla^\star f(x)}_2 \norm{\nabla^\star f(x^i)}_2 - \norm{\nabla^\star f(x)}_2^2}. \label{eq:norm-L-norm}
    \end{align}
    We then deduce from the Cauchy--Schwarz inequality that
    \begin{align}
        \norm{\nabla^\star f(x)}_2 \norm{\nabla^\star f(x^i)}_2
        &\ge \inner{\nabla^\star f(x)}{\nabla^\star f(x^i)} + 2|\delta^\star_i f(x)| |\delta^\star_i f(x^i)| \nonumber\\
        &= \inner{\nabla^\star f(x)}{\nabla^\star f(x^i)} + 2 \frac{q^\star_i(x^i)}{q^\star_i(x)} (\delta^\star_i f(x))^2. \label{eq:cauchy-schwarz}
    \end{align}
    where we crucially use the fact that $\delta^\star_i f(x) = q^\star_i(x) (f(x^i)-f(x))$ and $\delta^\star_i f(x^i) = q^\star_i(x^i) (f(x)-f(x^i))$ have opposite sign.
    Combining \Cref{eq:norm-L-norm,eq:cauchy-schwarz},
    \begin{align*}
        \norm{\nabla^\star f}_2 L \norm{\nabla^\star f}_2 (x)
        &\ge \sum_{i=1}^n q^\star_i(x) \tp{ \inner{\nabla^\star f(x)}{\nabla^\star f(x^i)} + 2 \frac{q^\star_i(x^i)}{q^\star_i(x)} (\delta^\star_i f(x))^2 - \inner{\nabla^\star f(x)}{\nabla^\star f(x)} } \\
        &= \sum_{i=1}^n 2 q^\star_i(x^i) (\delta^\star_i f(x))^2 + \sum_{i=1}^n q^\star_i(x) \inner{\nabla^\star f(x)}{\nabla^\star f(x^i) - \nabla^\star f(x)} \\
        &= \inner{\grad^\star f}{D \grad^\star f} (x) + \inner{\nabla^\star f}{L \nabla^\star f} (x).
    \end{align*}
    Plugging into \Cref{eq:first-rBE} and combining \Cref{lem:commutator}, we deduce
    \begin{align*}
        \Gamma^\star_2(f) - \Gamma(\sqrt{\Gamma^\star(f)}) 
        &\ge \inner{\grad^\star f}{D \grad^\star f} + \inner{\nabla^\star f}{(L \nabla^\star - \nabla^\star L) f} \\
        &= \inner{\grad^\star f}{\left( D + \diag\{M\*1\} - M^\top \right) \grad^\star f}
        \ge \rho \norm{\grad^\star f}_2^2 = \rho \cdot \Gamma^\star(f),
    \end{align*}
    where the last inequality holds for any constant $\rho \le \lambda_{\min} \left( D + \diag\{M\*1\} - \frac{1}{2}(M+M^\top) \right)$ uniformly at all points $x \in \Omega$.
\end{proof}

\subsection{Sufficient condition for CSGE} 
\label{sec:CSGE_boolean_hypercube}

\begin{theorem}[CSGE under canonical transition rates and weights]
\label{thm:CSGE_boolean_hypercube_w=q}
    Let $\mu$ be a distribution on $\Omega = \{\pm 1\}^n$ with full support. Let $M,D: \Omega \to \R^{n \times n}$ be two matrix fields defined by \Cref{eq:local-matrix,eq:def-D}, respectively. Let $M^\star := (\sup_x|M_{ij}(x)|)_{i,j \in [n]}$ and $D^\star = (\inf_x D_{ij}(x))_{i,j \in [n]}$ denote the entrywise uniform upper and lower bound on $\abs{M}$ and $D$, respectively. 
    Then, under the canonical transition rates and weights, for any $\tau \in \R_+$, \ref{eq:CSGE} holds with constant
    \begin{align*}
        \rho =
        \lambda_{\min}\left( D^\star - \diag\{M^\star\*1\} - \frac{1}{2}\left( M^\star + (M^\star)^\top \right) \right),
    \end{align*}
    where $\lambda_{\min}(A)$ denotes the minimum eigenvalue of a symmetric matrix $A$.
\end{theorem}

\ZC{How to deal with that $F_s$ is not differentiable at $s$?}

\ZJ{Does the following idea work?

$F_s$ is not differentiable at $s$ implies that there exists $i\in [n], x\in \Omega$, $\delta_i^\star P_{t-s}f(x) = 0$?
Fix $i$ and $x$, let $g(s) = P_{t-s}f(x)$.
% Since $g$ is in $C^\infty$ and $g_s(x) = \sum_{k=0}^\infty \frac{(t-s)^k}{k!}(L^k f)(x)$ converges, $g$ is an analytic function.
For any $s,s_0\in [0,t]$, $g(s) = \sum_{i=0}^{\infty} \frac{(L^i P_{t-s_0} f(x))}{i!} (s_0-s)^i$ converges.
Therefore, $g$ is an analytic function. (?)
Then $s\mapsto \delta_i^\star P_{t-s}f(x)$ is also an analytic function.
There are finitely many zeros for an analytic function $h$ in a closed interval unless $h\equiv 0$.
Therefore, $F_s$ is differentiable almost everywhere.
}

\begin{proof}
    Fix an arbitrary function $f: \Omega \to \R$ and real numbers $\tau, t \ge 0$.
    Our goal is to show \ref{eq:CSGE} under the canonical rates and weights, namely, 
    \begin{align}\label{eq:goal-CSGE}
        \norm{P_\tau \abs{\nabla^\star}P_t f}_2 \leq e^{-\rho t} \norm{P_{\tau+t} \abs{\nabla^\star}f}_2. 
    \end{align}
    Define $F_s = e^{-2\rho s} \norm{P_{\tau+s}\abs{\nabla^\star} P_{t-s} f}_2^2$; hence, \Cref{eq:goal-CSGE} is equivalent to $F_0 \le F_t$.
    It is sufficient to show that $\frac d{ds} F_s \geq 0$ for all $s \in [0,t]$.
    
    Fix $s \in [0,t]$. To ease the notation, let $g = P_{t-s} f$.
    Taking the derivative of $F_s$ gives 
    \begin{align}
        \frac d{ds} F_s &= 2 e^{-2\rho s} \left( - \rho \norm{P_{\tau+s} |\grad^\star| g}_2^2 
        + \inner{P_{\tau+s} |\grad^\star| g}{\frac{d}{ds} P_{\tau+s} |\grad^\star| g}
        \right).
        %\\&=2e^{-2\rho s}\inner{P_{\tau+s} |\grad^\star| g}{P_{\tau+s}\left(-\rho |\grad^\star| g + L |\grad^\star| g - \sgn(\grad^\star g)\otimes \grad^\star Lg\right)}.
        \label{eq:CSGE_potential_function_derivative}
    \end{align}
    Since $P_{\tau+s} = e^{(\tau+s)L}$ and $g = P_{t-s} f = e^{(t-s)L} f$, we have that
    \begin{align}
    \label{eq:derivative-CSGE}
        \frac d{ds} P_{\tau+s} |\grad^\star| g
        &= P_{\tau+s} \left( L + \frac{d}{ds} \right) |\grad^\star| g 
        = P_{\tau+s} \left( L |\grad^\star| - \sgn(\grad^\star g) \odot \grad^\star L \right) g,
    \end{align}
    where $\sgn(\cdot)$ denotes the entrywise absolute value of a vector, and $\odot$ denotes the entrywise product of two vectors.
    To calculate this, we need an adapted version of \Cref{lem:commutator}, which is stated below and proved afterward.

    \begin{lemma}
    \label{lem:abs-commutator}
        Let $\mu$ be a distribution supported on $\Omega = \{\pm 1\}^n$.
        Let $M,D: \Omega \to \R^{n \times n}$ be two matrix fields defined by \Cref{eq:local-matrix,eq:def-D}, respectively, and let $|M| := (|M_{ij}|)_{i,j \in [n]}$ be the entrywise absolute value of $M$.
        For any function $g: \Omega \to \R$, it holds
        \begin{align*}
            \left( L |\grad^\star| - \sgn(\grad^\star g) \odot \grad^\star L \right) g \ge \left( D + \diag\{M\*1\} - |M|^\top \right) |\grad^\star| g.
        \end{align*}
        % \ZC{I got
        % \begin{align*}
        %     \left( L |\grad^\star| - \sgn(\grad^\star g) \odot \grad^\star L \right) g \ge \left( D + \diag\{M\*1\} - |M|^\top \right) |\grad^\star| g,
        % \end{align*}
        % }
        % \ZJ{We can use this tighter version, although I don't think it can improve the final conclusion.} \ZC{Agree!}
    \end{lemma}
    
    Combining \Cref{eq:derivative-CSGE,lem:abs-commutator}, 
    we deduce that
    \begin{align*}
        \inner{P_{\tau+s} |\grad^\star| g}{\frac{d}{ds} P_{\tau+s} |\grad^\star| g}
        &\geq \inner{P_{\tau+s} |\grad^\star| g}{P_{\tau+s} \left( D + \diag\{M\*1\} - |M|^\top \right) |\grad^\star| g}
        \\&\geq \inner{P_{\tau+s} |\grad^\star| g}{ \left( D^\star - \diag\{M^\star\*1\} -( M^\star)^\top \right) P_{\tau+s}|\grad^\star| g}.
    \end{align*}
    % \ZC{
    % Call $A := D + \diag\{M\*1\} - |M|^\top$, which is a matrix field $A: \Omega \to \R^{n \times n}$.
    % It is an operator acting on a vector field $v: \Omega \to \R^n$ via pointwise product:
    % \begin{align*}
    %     (Av)(x) := A(x) v(x), \quad\forall x \in \Omega.
    % \end{align*}
    % Meanwhile, $P_{\tau+s}$ is another operator acting on vector fields entrywisely:
    % \begin{align*}
    %     P_{\tau+s} v := (P_{\tau+s} v_1, \dots, P_{\tau+s} v_n).
    % \end{align*}
    % Your proof seems to claim they commute:
    % \begin{align*}
    %     P_{\tau+s} A v = A P_{\tau+s} v.
    % \end{align*}
    % Do you have a proof? This feels suspicious.\\
    % A possible fix is to define $A$ uniformly:
    % \begin{align*}
    %     A_{ij} := 
    %     \begin{cases}
    %         \min_{x \in \Omega} D_{ii}(x) + \sum_{j=1}^n M_{ij}(x), & i = j \\
    %         - \max_{x \in \Omega} |M_{ji}(x)|, & i \neq j 
    %     \end{cases}
    % \end{align*}
    % I would keep the statement in \Cref{lem:abs-commutator} the same except adding
    % \begin{align*}
    %         \left( L |\grad^\star| - \sgn(\grad^\star g) \odot \grad^\star L \right) g 
    %         \ge \left( D + \diag\{M\*1\} - |M|^\top \right) |\grad^\star| g 
    %         \ge A |\grad^\star| g.
    %     \end{align*}
    % }
    We conclude that 
    $$
        \frac{d}{ds}F_s \geq \tp{P_{\tau+s} |\grad^\star| g}^\top\tp{D^\star - \diag\{M^\star\*1\}-\frac{1}{2}(M^\star+(M^\star)^\top) - \rho I}\tp{P_{\tau+s} |\grad^\star| g} .
    $$
    Therefore, the coordinate-wise strong gradient estimate (\ref{eq:CSGE}) holds with constant $\rho = \lambda_{\min}(D^\star-\diag\{M^\star\*1\} - \frac{1}{2}(M^\star+(M^\star)^\top))$ at all $x \in \Omega$.
\end{proof}

\begin{proof}[Proof of \Cref{lem:abs-commutator}]
    % \ZC{Update the proof.}
    Define $H_{ij}(x) = \frac{q_i(x^j)}{q_i(x)}$.
    By \Cref{fact:square}, we have that $H_{ij}(x) = \frac{q_i(x^j)}{q_i(x)} = \frac{q_j(x^i)}{q_j(x)} = H_{ji}(x)$.
    For notational convenience, we occasionally omit the dependence on $x$, writing, for example, $q_i^\star$ instead of $q_i^\star(x)$.
    For $i,j\in[n]$ with $i\neq j$, the following is deduced from $H_{ij}=H_{ji}$:
    \begin{align}
        (\delta_j^\star \abs{\delta_i^\star} - \sgn(\delta^\star_i g(x))\delta_i^\star \delta_j^\star) g(x)
        &=  q_j^\star q_i^\star\tp{H_{ij}\abs{g(x^{ij}) - g(x^j)} - \abs{g(x^i) - g(x)}} \nonumber
        \\&- \sgn(\delta^\star_i g) q_i^\star q_j^\star(H_{ji}(g(x^{ij})-g(x^i)) - g(x^j)+g(x))\nonumber
        \\&=q_i^\star q_j^\star H_{ij}\tp{\abs{g(x^{ij}) - g(x^j)} - \sgn(\delta^\star_i g)(g(x^{ij}) - g(x^j))} \label{eq:abs-commutator_second_term}
        \\&+q_i^\star q_j^\star\tp{-\frac{\abs{\delta^\star_i} g}{q_i^\star} - \sgn(\delta^\star_i g)\tp{H_{ij}(g(x^j) - g(x^i)) - \frac{\delta^\star_jg}{q_j^\star}}}. \label{eq:abs-commutator_eq}
    \end{align}
    Since $\abs{a}-\sgn(b)a \geq 0$ holds for any $a,b\in\R$, \Cref{eq:abs-commutator_second_term} is non-negative, resulting in the callcelation of the second derivative term in the commutator.
    Therefore, by \Cref{eq:abs-commutator_eq}, 
    \begin{align}
        (\delta_j^\star \abs{\delta_i^\star} - \sgn(\delta^\star_i g(x))\delta_i^\star \delta_j^\star) g(x)
        &\geq - q_j^\star\abs{\delta^\star_i g} - \sgn(\delta^\star_i g)\tp{H_{ij}q_i^\star q_j^\star(\frac{\delta_j^\star g}{q_j^\star} - \frac{\delta_i^\star g}{q_i^\star}) - q_i^\star\delta^\star_jg} \nonumber
        \\&= (H_{ij} - 1)q_j^\star\abs{\delta^\star_i}g - \sgn(\delta^\star_i g) \tp{H_{ij} - 1}q_i^\star\delta^\star_j g \nonumber
        \\&\geq (H_{ij} - 1)q_j^\star\abs{\delta^\star_i}g - \abs{H_{ji} - 1}q_i^\star\abs{\delta^\star_j}g \label{eq:CSGE_commutator_term}
    \end{align}
    For the non-commutator term, we have that
    \begin{align} \label{eq:CSGE_non-commutator_term}
        (\delta_i^\star \abs{\delta_i^\star} - \sgn(\delta^\star_i g(x))\delta_i^\star \delta_i^\star)g(x) &= q_i^\star(x) \tp{\abs{\delta_i^\star}g(x^i) - \sgn(\delta^\star_i g(x))\delta_i^\star g(x^i)} \nonumber
        \\&= 2 q_i^\star(x) \abs{\delta_i^\star}g(x^i)
        = 2 q_i^\star(x^i) \abs{\delta_i^\star}g(x)
        = \frac 1{2q_i^\star(x)} \abs{\delta_i^\star}g(x),
    \end{align}
    where we use the fact that for any $i\in[n]$, $x^{ii} = x$ and $\sgn(\delta^\star_i g(x)) = \sgn(g(x^i) - g(x)) = - \sgn(g(x) - g(x^i)) = -\sgn(\delta^\star_i g(x^i))$.
    Combine \Cref{eq:CSGE_commutator_term,eq:CSGE_non-commutator_term} together, we conclude the following inequality,
    \begin{align*}
        (L\abs{\delta_i^\star} - \sgn(\delta^\star_i g(x))\delta_i^\star L) g(x) &\geq \frac{1}{2q_i^\star}\abs{\delta^\star_i} g + \sum_{j\neq i} (H_{ij} - 1)q_j^\star\abs{\delta^\star_i} g - \abs{H_{ji} - 1}q_i^\star\abs{\delta^\star_j} g
        \\&= (D_{ii} + (M_i)^\top \*1)\abs{\delta^\star_i} g - \inner{(\abs{M}^\top)_i}{\abs{\nabla^\star} g}, 
    \end{align*}
    where the last equality follows from $\abs{H_{ij}-1}p_j^\star = \abs{M_{ij}}$.
    This is exactly
    \[
        \left( L |\grad^\star| - \sgn(\grad^\star g) \odot \grad^\star L \right) g \ge \left( D + \diag\{M\*1\} - |M|^\top \right) |\grad^\star| g. \qedhere
    \]
\end{proof}


\subsection{Dobrushin condition}
In this section, we state our isoperimetric inequalities under the conditions of the Dobrushin matrix and marginal bounds. 
In addition to the CSGE condition that we establish in the above, we also require that hypercontractivity holds from \Cref{sec:boolean}.

It turns out that the largest eigenvalue of $D$ and $M$ can be controlled by suitable norms of Dobrushin influence matrix with some additional bounds associated with the canonical transition rates.
Consequently, a general and easy way to check the CSGE condition is to establish an upper bound the norms of the Dobrushin matrix and check the marginal bounds of the distribution $\mu$.
The following is the definition of the Dobrushin matrix.
\begin{definition}[Dobrushin influence]
    For finite space $\+X$, the Dobruhsin influence matrix of distribution $\mu$ over $\+X^n$ is defined as follows,
    \begin{align}\label{eq:def_Dobrushin_matrix}
        R = (R_{ij})_{i,j\in[n]},\quad R_{ij} = \sup_{x,y,x_{-j}=y_{-j}}\TV{\mu_i(\cdot\mid x_{-i})}{\mu_i(\cdot\mid y_{-i})} .
    \end{align}
    We say a distribution satisfies the Dobrushin condition if $\norm{R(\mu)}_2 < 1$.
    \ZC{Are we defining $R$ in the wrong way? $R \gets R^\top$}
    \ZJ{I saw different definitions, the current definition is the same as the definition from Frederic's \href{chrome-extension://efaidnbmnnnibpcajpcglclefindmkaj/https://frkoehle.github.io/stat31512-s2024/lecture10.pdf}{lecture notes}.}
\end{definition}

\begin{lemma}[]\label{lem:Dobrushin_upper_bound_CSGE}
     If $w_i^\star(x) \in [b,1/b]$ for any $x\in \{\pm 1\}^n$ and $i\in[n]$, then under the canonical transition rates and weights, \ref{eq:CSGE} holds with constant $\rho = b - \frac{(1+b^2)^2}{4b^3}(\norm{R}_2 + \norm{R}_1)$.
\end{lemma}
\begin{proof}
    The entries of the Dobrushin influence matrix can be expressed in terms of the canonical transition rates as follows:
    $$
        R_{ij} = \begin{cases}
            \sup_x \frac{4\abs{q_i^\star(x^j)^2-q_i^\star(x)^2}}{(1+4q_i^\star(x)^2)(1 + 4q_i^\star(x^j)^2)} &, i\neq j
            \\ 0 &, i = j
        \end{cases} .
    $$
    Then the Dobrushin matrix give an entry-wise upper bound of $M^\star = (\sup_x \abs{M_{ij}(x)})_{i,j\in[n]}$.
    That is, for any $i, j\in[n]$ and $i\neq j$,
    \begin{align}\label{eq:Dobrushin_to_CSGE_constant}
        R_{ij} \geq \tp{\inf_{u,v\in [b,1/b]}\frac{u+v}{(1+u^2)(1+v^2)}}\sup_{x\in \Omega}2\abs{q_i^\star(x)-q_i^\star(x^j)} \geq \frac{4b^3}{(1+b^2)^2}M^\star_{ji} .
    \end{align}
    % \ZC{Our goal is to show: 
    % \begin{align}
    %     A - \diag\{B\*1\}-\rho I- \frac{B+B^{\top}}{2}\succeq 0.
    % \end{align}
    % Why do we need $K$?
    % I think we have 
    % \begin{align}
    %     A \succeq c_1 I, \quad \diag\{B\*1\} \preceq c_2 I, \quad B+B^{\top} \preceq c_3 \lambda_{\max}(D+D^\top) \cdot I.
    % \end{align}
    % }
    Since $D^\star_{ii} = \inf_x w_i(x^i)$ for $i\in[n]$, $D^\star\geq_{\-{ew}} bI$.
    Therefore, by symmetry of matrices and \Cref{eq:Dobrushin_to_CSGE_constant}, we have that
    \begin{align*}
        \lambda_{\min}(D^\star-\diag\{M^\star\*1\}-\frac{1}{2}(M^\star+(M^\star)^\top)) &\geq \lambda_{\min}(D^\star) - \lambda_{\max}(\diag\{M^\star\*1\}) - \lambda_{\max}(\frac{1}{2}(M^\star+(M^\star)^\top)) 
        \\&\geq b - \frac{(1+b^2)^2}{4b^3}(\norm{R}_2 + \norm{R}_1)  \geq \rho.
    \end{align*}
    By \Cref{thm:CSGE_boolean_hypercube_w=q}, we complete the proof.
\end{proof}

The following theorem shows that for the discrete product space, the Dobrushin condition and the one-site conditional marginal lower bound are sufficient to derive \ref{eq:LSI} constant and hypercontractivity. 
\begin{theorem}[{\cite[Theorem 1.14 and Corollary 1.11]{marton2019logarithmic}}]\label{thm:marton_Dobrushin_to_LSI}
Let $\+X$ be a finite space.
Suppose the distribution $\mu$ on $\Omega = \+X^n$ satisfies the Dobrushin condition $\norm{R}_2 < 1$ and the one-site conditional marginal lower bound $\alpha = \inf_x \mu(x_i\mid x_{-i})$, then for Gibbs sampler $G = \frac 1 n\sum_{i=1}^n K_i$ where $K_i(z\mid y) = \ind[z_{-i}=y_{-i}]\mu(z_i\mid z_{-i})$, the log-Sobolev inequality holds as follows,
$$
    \forall f\in \R^\Omega, \frac 1n\Ent[\mu]{f^2} \leq \frac{2}{\alpha(1-\norm{R}_2)^2} \inner{f}{(I-G)f}_{\mu} .
$$
\end{theorem}
\begin{corollary}\label{cor:Dobrushin_to_LSI}
    Suppose $w_i^\star(x)\in [b,1/b]$ for any $x\in \{\pm 1\}^n$ and $i\in[n]$, then \ref{eq:LSI} holds with constant $\frac{b^2(1 - \norm{R}_2)^2}{1+b^2}$ with respect to the canonical transition rates $q^\star$.
\end{corollary}
\begin{proof}
    This immediately follows from
    \[
        \inner{f}{(I-G)f}_{\mu} = \E[\mu]{\frac 12\sum_{i=1}^n \frac{2q_i^\star(x)}{n(1/w_i^\star(x)+w_i^\star(x))}(f(x^i)-f(x))^2}
        \leq \frac 1n\E[\mu]{\Gamma(f,f)} . \qedhere
    \]
\end{proof}
By \Cref{thm:CSGE_boolean_hypercube_w=q}, \Cref{cor:Dobrushin_to_LSI} and \cref{sec:boolean}, we conclude the following theorem about isoperimetric inequalities.
\begin{theorem}\label{thm:Dobrushin_to_isoperimetric}
    Let $\mu$ be a distribution on $\{\pm 1\}^n$ with full support.
    Let the transition rates and weights to be the canonical ones defined in \Cref{eq:canonical-r-w}.
    Suppose that for any $x\in \{\pm 1\}^n$ and $i\in[n]$, $w_i^\star(x)\in [b,1/b]$ and $\max\{\norm{R}_1,\norm{R}_2\}\leq C$ for some constants $b,C$, where $R$ is the Dobrushin influence matrix defined in \Cref{eq:def_Dobrushin_matrix}.
    Then the local Bobkov inequality \eqref{eq:local_bobkov} holds.

    Furthermore, if the Dobrushin condition holds, the $L^1$--$L^2$ Talagrand inequality \eqref{eq:l1_l2_talagrand_2}, the KKL inequality \eqref{eq:} and the Eldan-Gross inequality \eqref{eq:eldan_gross} holds with constant depending on $\norm{R}_1$, $\norm{R}_2$ and $b$.
\end{theorem}
\subsection{Application to Ising model}
A distribution $\mu$ over $\{\pm 1\}^n$ is called an \emph{Ising model} with interaction matrix $J \in \R^{n\times n}$ and external field $h\in \R^n$ if 
$$
    \forall x\in \{\pm 1\}^n, \mu(x) \propto \exp(\frac 12 x^\top J x + h^\top x),
$$
where $J$ is a symmetric matrix.
Note that the diagonals of $J$ do not affect the distribution, hence we can assume that $J_{ii} = 0$ for $i\in[n]$.
\ZJ{previous results.}

In the following theorem, we establish a new criterion for high temperature Ising model to prove the isoperimetric inequalities under a Dobrushin-type condition by \cref{thm:Dobrushin_to_isoperimetric}.
\begin{theorem}\label{thm:Ising}
    Suppose that $\mu$ is an Ising model parameterized by zero external field $h = \*0$ and interaction matrix $J$ satisfying $\norm{J}_2 < 1$ and $\norm{J}_1\leq C$ for some constant $C > 0$.
    Let the transition rates and weights to be the canonical ones defined in \Cref{eq:canonical-r-w}.
    Then the local Bobkov inequality \eqref{eq:local_bobkov}, the $L^1$--$L^2$ Talagrand inequality \eqref{eq:l1_l2_talagrand_2}, the KKL inequality \eqref{eq:} and the Eldan-Gross inequality \eqref{eq:eldan_gross} holds with constant depending on $\norm{J}_2$ and $C$.
\end{theorem}


\begin{proof}
    Let $b = e^{-\norm{J}_1}$.
    For any $i,x$, $w_i^\star(x) = \exp(-x_iJ_i^\top x) \in [b,1/b]$.
    Fix $x\in \Omega$ and $i,j\in [n]$ with $i\neq j$.
    Define $\alpha = \sum_{k\neq i,j}J_{ik} x_k$ and $\beta = \abs{J_{ij}}$ .
    We have that 
    \begin{align*}
        \TV{\mu_i(\cdot\mid x)}{\mu_i(\cdot\mid x^j)} &= \abs{\mu(x_i=+1\mid x_j=+1,x_{-i,j}) - \mu(x_i=+1\mid x_j=-1,x_{-i,j})}
        \\&= \frac 12\abs{\tanh(\alpha+\beta) - \tanh(\alpha-\beta)}
        = \frac{2\abs{\sinh(\beta)\cosh(\beta)}}{\cosh(2\alpha)+\cosh(2\beta)}
        \\&\leq \frac{2\abs{\sinh(\beta)\cosh(\beta)}}{1+\cosh(2\beta)} = \tanh(\beta) \leq \beta.
    \end{align*}
    Therefore, $R \leq_{\-{ew}} \abs{J}$ and $\norm{R}_2\leq \norm{\abs{J}}_2 < 1$.
    We have that $\norm{J}_2 \leq \sqrt{\norm{J}_1\norm{J}_{\infty}} = \norm{J}_1$ follows from symmetry. 
    By \cref{thm:Dobrushin_to_isoperimetric}, since $\max\{\norm{R}_1,\norm{R}_2\}\leq \norm{J}_1 \leq C$, all the isopermetric inequalities in \cref{sec:boolean} hold.
\end{proof}



